% =====================================================================
%  Quantitative retrieval evaluation.
%  Insert after Section 3.3 (Qualitative Evaluation and Comparison).
%
%  Every number here is measured. Reproduce with:
%      uv run python eval/run_eval.py --configs bm25,vector,rrf,quality,full
%      uv run python eval/run_eval.py --queries eval/queries_topical.jsonl ...
% =====================================================================

\section{Quantitative Retrieval Evaluation}
\label{sec:retrieval-eval}

The evaluation in Section~\ref{sec:qualitative-eval} scores complete
recommendations with an LLM judge. That measures what a user finally sees, but
it cannot say \emph{where} a poor result came from. A single composite score
does not distinguish a game that was never retrieved from one that was
retrieved and then ranked too low, and those two faults are repaired in
different parts of the pipeline. Five prompts also leave a mean whose
confidence interval is wide enough that a ten-point swing on one prompt moves
it by two.

This section adds an automatic evaluation of the retrieval and ranking stages.
It complements rather than replaces the judge protocol: it says nothing about
the generated text, and the judge says nothing about which stage failed.

\subsection{Two Sources of Ground Truth}
\label{sec:eval-tracks}

Neither available source of relevance judgements is sound alone, so both are
used and read against each other.

\paragraph{Track A --- known-item retrieval.}
A held-out player review is used as the query, and the game it was written
about is the correct answer. No annotation is required and the label cannot be
disputed. Validity rests on two properties. The review must not already be
indexed: the index embeds, per game, the fifteen most-upvoted positive reviews
within a length window, so queries are drawn at rank sixteen under the
identical ordering. And titles must be removed, because 30\% of reviews name
their own game, which lexical retrieval would otherwise match directly. 500
queries were generated this way; 26,277 of the 30,693 indexed games have a
review available beyond what is indexed.

\paragraph{Track B --- topical relevance.}
Fifty hand-written queries, ten for each of the five intent categories used in
Section~\ref{sec:qualitative-eval}, graded against tag sets: grade 2 when every
required tag is present, grade 1 for a partial match, 0 otherwise. Hard
constraints on price, platform and multiplayer support, and exclusions for the
negation category, force a grade of 0 when unmet. Query wording deliberately
avoids the tag strings themselves, so what is measured is the ability to bridge
a vocabulary gap rather than to match a literal.

\paragraph{Why both.}
Track B is closer to the real task but partly circular: tags are part of the
indexed corpus as well as the source of the labels, which favours lexical
matching. Track A is free of that circularity but asks a different question ---
find one specific game, rather than assemble a useful set. As
Section~\ref{sec:eval-attribution} shows, the two tracks disagree about which
components matter, and the disagreement is the most informative result here.

\subsection{Results}
\label{sec:eval-results}

\begin{table}[h]
\centering
\caption{Track A, known-item retrieval over 500 queries. Each query has exactly
one correct answer, so P@5 cannot exceed 0.200. Random selection would give
Recall@50 $\approx$ 0.0016. Values are mean $\pm$ standard error.}
\label{tab:track-a}
\small
\begin{tabular}{llcccc}
\hline
\textbf{\#} & \textbf{Configuration} & \textbf{Recall@50} & \textbf{nDCG@10} & \textbf{MRR} & \textbf{P@5} \\
\hline
1 & Dense, 18k index (no prefix) & 0.134 $\pm$ 0.015 & 0.043 $\pm$ 0.008 & 0.037 $\pm$ 0.007 & 0.010 $\pm$ 0.002 \\
2 & BM25                         & 0.180 $\pm$ 0.017 & 0.073 $\pm$ 0.010 & 0.065 $\pm$ 0.009 & 0.018 $\pm$ 0.003 \\
3 & Dense, 30k index             & 0.182 $\pm$ 0.017 & 0.070 $\pm$ 0.010 & 0.059 $\pm$ 0.009 & 0.018 $\pm$ 0.003 \\
4 & \textbf{Rank fusion (2 + 3)} & \textbf{0.248 $\pm$ 0.019} & \textbf{0.102 $\pm$ 0.011} & \textbf{0.086 $\pm$ 0.010} & 0.024 $\pm$ 0.003 \\
5 & \quad + quality \& diversity & 0.248 $\pm$ 0.019 & 0.098 $\pm$ 0.011 & 0.084 $\pm$ 0.010 & 0.026 $\pm$ 0.003 \\
\hline
\end{tabular}
\end{table}

\begin{table}[h]
\centering
\caption{Track B, topical relevance over 50 queries. A median of 1,206 games
score grade 2 per query, so Recall@50 cannot exceed about 0.02 and is not
reported. The median random baseline for P@5 is 0.039.}
\label{tab:track-b}
\small
\begin{tabular}{llcccc}
\hline
\textbf{\#} & \textbf{Configuration} & \textbf{nDCG@10} & \textbf{MRR} & \textbf{P@5} & \textbf{ILD@5} \\
\hline
1 & Dense, 18k index (no prefix) & 0.345 $\pm$ 0.043 & 0.597 $\pm$ 0.058 & 0.408 $\pm$ 0.047 & --- \\
2 & BM25                         & 0.341 $\pm$ 0.033 & 0.588 $\pm$ 0.053 & 0.408 $\pm$ 0.041 & --- \\
3 & Dense, 30k index             & 0.356 $\pm$ 0.044 & 0.562 $\pm$ 0.057 & 0.404 $\pm$ 0.048 & --- \\
4 & Rank fusion (2 + 3)          & 0.389 $\pm$ 0.039 & 0.634 $\pm$ 0.053 & 0.420 $\pm$ 0.040 & 0.881 $\pm$ 0.009 \\
5 & \quad + review-aware quality & 0.429 $\pm$ 0.040 & 0.724 $\pm$ 0.050 & 0.480 $\pm$ 0.043 & 0.872 $\pm$ 0.011 \\
6 & \quad\quad + diversity       & \textbf{0.440 $\pm$ 0.040} & \textbf{0.751 $\pm$ 0.050} & \textbf{0.496 $\pm$ 0.043} & 0.873 $\pm$ 0.011 \\
\hline
\end{tabular}
\end{table}

The Recall@50 figures in rows 4 and 5 of Table~\ref{tab:track-a} are identical
by construction: re-ranking reorders a candidate set without changing its
membership. A difference there would indicate a fault in the harness rather
than an effect of the component, so the equality is a check on the experiment.

\subsection{Component Attribution}
\label{sec:eval-attribution}

Every configuration was run over the same queries, so differences are paired
and their standard errors are considerably tighter than those in the tables
above.

\begin{table}[h]
\centering
\caption{Paired differences between successive configurations. Significance is
the ratio of the paired mean difference to its standard error.}
\label{tab:attribution}
\small
\begin{tabular}{lllc}
\hline
\textbf{Component} & \textbf{Track} & \textbf{Effect} & \textbf{Verdict} \\
\hline
Rank fusion
  & A & Recall@50 $+0.066$ ($3.5\sigma$) & confirmed \\
  & B & nDCG $+0.033$ ($1.1\sigma$), P@5 $+0.016$ ($0.4\sigma$) & trend only \\[2pt]
Review-aware quality
  & B & nDCG $+0.039$ ($3.6\sigma$), MRR $+0.090$ ($2.9\sigma$), P@5 $+0.060$ ($3.5\sigma$) & confirmed \\
  & A & nDCG $-0.004$ & not measurable \\[2pt]
Diversity constraint
  & B & nDCG $+0.011$ ($0.8\sigma$), P@5 $+0.016$ ($1.1\sigma$), ILD $+0.001$ & not confirmed \\[2pt]
Corpus coverage
  & A & Recall@50 $+0.048$ & confirmed \\[2pt]
Task prefixes, review sampling
  & A & Recall@50 $+0.000$ on the common subset & not confirmed \\
  & B & nDCG $+0.011$ ($0.4\sigma$), MRR $-0.035$ ($-0.7\sigma$) & not confirmed \\[2pt]
Reviews in the lexical corpus
  & A & Recall@50 $0.000 \rightarrow 0.133$ & confirmed \\
\hline
\end{tabular}
\end{table}

\paragraph{The two tracks disagree, and both are right.}
Rank fusion is decisive on Track A and marginal on Track B; review-aware
quality scoring is the reverse. Once the tasks are distinguished this is what
should be expected. Fusion pays off when two retrievers fail on different
queries, and a topical query with over a thousand valid answers is one that
either retriever already partly satisfies. Quality scoring reorders candidates
that are all plausible, which is the situation on Track B and not on Track A,
where 499 of the 500 candidates are simply wrong. Running one track alone would
have produced a wrong conclusion in either direction.

\paragraph{Fusion combines retrievers that score alike but differ.}
On Track A, BM25 and dense retrieval reach 0.180 and 0.182 --- statistically
indistinguishable --- yet their fusion reaches 0.248, a 36\% relative gain.
Equal aggregate scores with a large gain from combining them is direct evidence
that the two fail on \emph{different} queries, which is the premise the hybrid
design rests on and which the individual scores alone cannot establish.

\paragraph{Corpus content matters more than corpus selection.}
The largest single effect in the table comes from adding reviews to the lexical
corpus. Before that change the BM25 configuration scored 0.000 on every Track A
metric; a control test confirmed the retriever worked, ranking the target first
when queried by title or store description. The queries are player reviews, and
the lexical corpus contained only marketing copy and tag lists, so there was
almost no shared vocabulary. By contrast, changing \emph{which} fifteen reviews
the dense index embeds --- from the longest to the most upvoted --- produced no
measurable effect on either track. Whether a field is present matters far more
than which part of it is selected.

\paragraph{A mechanically sound fix with no measurable benefit.}
The embedding model is asymmetric and documents a required task prefix for
queries and documents. Neither side applied one. The repair was implemented and
the index rebuilt, and neither track detects a gain; Track B's MRR is
nominally negative. Isolating it required care, because 147 of the 500 Track A
gold games are absent from the old index and cannot be found there at all. On
the 353 queries answerable by both, Recall@50 is 0.190 either way, which
attributes the entire corpus gain to coverage and none to the prefixes.

We report this as a negative result rather than omitting it. That a model's
documentation prescribes a prefix does not establish that adding it improves
retrieval on a particular dataset and task; without this evaluation the change
would have been described in this report as a fixed defect.

\paragraph{A constraint that rarely binds.}
The diversity pass caps how many results may share a developer or a primary
tag. It changes the top five on 7 of 50 queries and the top result on 2, and
intra-list diversity is unchanged at 0.87 with or without it. The fused
shortlists are already varied, so the caps seldom apply. The failure mode the
constraint was written for --- five entries of one series filling the
shortlist --- is real in principle but rare enough here that the mechanism
cannot be shown to help.

\subsection{Limitations}
\label{sec:eval-limitations}

\paragraph{Tag labels are partly circular.}
Track B grades relevance by tags, and tags are also part of the indexed corpus,
which favours lexical retrieval. Track A exists partly to provide a check free
of this, and the two tracks agree on the direction of every confirmed
component.

\paragraph{Track A is a proxy for the real task.}
Finding the one game a review was written about is not the same as assembling
five good recommendations. Its value is that it needs no annotation and scales
to 500 queries; its limitation is that it rewards precision on a single target.

\paragraph{Absolute levels are low on Track A.}
Recall@50 of 0.248 means three quarters of queries fail. Locating one game
among 30,693 from a paraphrased review is hard --- random selection would score
0.0016, so the system is roughly 150 times better than chance --- and many
reviews contain nothing distinguishing (``great game, worth it''), which no
retrieval method can resolve. MRR divided by Recall@50 is 0.35, so when the
answer is retrieved at all it sits near rank three: recall, not ranking, is the
bottleneck.

\paragraph{Fifty queries limit Track B.}
Differences below roughly $2\sigma$ cannot be separated from noise at this
sample size, so ``not confirmed'' above means exactly that and not
``ineffective''.
